\documentclass[11pt,twoside]{article}
 
\usepackage{hyperref}
\usepackage{graphicx}
\usepackage{amsmath}

\pagestyle{plain}

\setlength{\voffset}{0pt}
\setlength{\marginparwidth}{0pt}
\setlength{\oddsidemargin}{30pt}
\setlength{\evensidemargin}{30pt}
\setlength{\textwidth}{400pt}
\setlength{\parindent}{0pt}
\setlength{\parskip}{10pt plus 2pt minus 1pt}

\title{\textbf{DruL} Reference Manual\\
\vspace{1cm}
COMS W4115: Programming Language and Translators}
\author{Team Leader: Rob Stewart (rs2660) \and Thierry Bertin-Mahieux (tb2332) \and Benjamin Warfield (bbw2108) \and Waseem Ilahi (wki2001)}

\begin{document}
\maketitle
\begin{center}
\end{center}

\vspace{3cm}

\begin{figure}[h]
\begin{center}
\includegraphics[width=.2\columnwidth]{Water_Drop.pdf}
\end{center}
\end{figure}


\newpage


\section{Introduction}
Here is the language manual.


\section{Language Manual}

\subsection{Lexical Conventions}

\subsubsection{Comments}
Comments start with // and everything after them on that line
is considered a comment and is ignored. New line marks the end
of a comment.
\subsubsection{Whitespace}
Spaces, tab, end of line, and return are all considered the same 
and their only used is to seperate the tokens.

\subsubsection{Characters}
We use ASCII characters.

\subsubsection{Keywords}


\begin{table}[htb]
\begin{center} 
\begin{tabular}{llll} 
  NULL &    rand &    clip &        mapper \\
  if &      pattern & instrument &  print \\
  elseif &  concat &  length &      output \\
  else &    slice &   map &         return
\end{tabular}
%\caption{\small{.}}
\label{tab:keywords}
\end{center} 
\end{table}


\subsubsection{Identifiers}
An identifier may be defined using any uppercase or lowercase character
or an underscore, followed by any set of uppercase or lowercase character,
underscore, and digits ($0$ through $9$). The maximum length of an identifier
is $64$.

\subsection{Types}

There are $3$ types in Drul: \textbf{int}, \textbf{pattern},
and \textbf{clip}.

\subsubsection{int}
We accept only positive integers, so any combination of digits
($0$ through $9$), such as $0000123$,  is interpreted as $123$.

\subsubsection{pattern}
A pattern is a series of events and silences.


\subsubsection{clip}
A clip is a set of patterns associated with instruments.

\subsection{Statements}

\subsubsection{Expressions}


\subsubsection{Blocks}
The curly braces mark the beginning and end of a block.

\subsubsection{Symbols}
$\&\&$, $||$, $==$, etc


\subsubsection{Logical expressions}
Do we have them?


\subsubsection{Conditional}


\subsubsection{Assignments}
We can assign any type to any identifier using the equal ($=$) sign.
???

\subsection{Patterns and pattern operations}

\subsubsection{Patterns}

\subsubsection{Map}

\subsubsection{Mapper}
Mapper has to be declared before it is used in the file.

\subsection{Clips}

\subsubsection{Instruments}


\subsubsection{Clips}

\subsection{Outputs}

\subsubsection{Standard output}

\subsubsection{Text}

\subsubsection{MIDI}

\end{document}
